% grrt RNAAS note (Phase 3D.3). Build: `make paper` (needs latexmk,
% pdflatex, or tectonic). Body word budget <= 1000: see `make wordcount`.
% AUTHOR/AFFILIATION/ORCID below are placeholders -- fill before submission.
\documentclass[RNAAS]{aastex631}

\begin{document}

\title{Photon-Orbit Bifurcation and ISCO Disappearance in a
Johannsen--Psaltis Ray Tracer at M87* Spin}

% TODO(author): replace the placeholder ORCID below before submission.
\author[0000-0000-0000-0000]{Pranav Sethuraman}
\affiliation{Independent Researcher, Fremont, CA, USA}

\begin{abstract}
We present \texttt{grrt}, an open-source general-relativistic ray tracer,
and use it to sweep the single Johannsen--Psaltis (JP) deviation parameter
$\epsilon_3$ at spin $a=0.9$; the code reproduces Kerr to machine
precision. We report two future-observable predictions---a prograde
photon-orbit bifurcation at $\epsilon_3^{\rm crit}=0.1212$ and prograde
ISCO disappearance for $\epsilon_3\in(0.13,0.20)$---and a monotonic
ring-circularity trend. Naive integer-pixel ring extraction floors at
$\sim\!12\%$, well above the EHT M87* circularity ($\sim\!5.5\%$); we
therefore claim no $\epsilon_3$ bound and identify sub-pixel extraction
as the required next step.
\end{abstract}

\section{Introduction}
The Event Horizon Telescope (EHT) image of M87* \citep{eht2019:paper-i}
is consistent with the shadow of a Kerr black hole, yet the data also
admit non-Kerr spacetimes. The Johannsen--Psaltis (JP) metric
\citep{johannsen:2011} parameterizes such deviations through a series
whose leading free term $\epsilon_3$ preserves the asymptotic and
weak-field limits. Shadow circularity is a natural diagnostic: the EHT
reports the M87* image circular to a fractional diameter spread of
$\sim\!5$--$6\%$ \citep[][\S7.4]{eht2019:paper-vi}. We present an
independent, open ray tracer and a controlled $\epsilon_3$ sweep at the
M87*-favored spin $a=0.9$, and ask what the sweep predicts and what
measurement precision a circularity test of $\epsilon_3$ would demand.
We do not claim a bound; we report the pipeline, two qualitative
predictions, and a methodological limit.

\section{Methods}
\texttt{grrt} is a Java~21 ray tracer (signature $-\!+\!+\!+$,
$G=c=1$). Each metric supplies closed-form, hand-derived Christoffel
symbols, verified symbolically against an independent \textsc{SymPy}
derivation; null geodesics are integrated with an adaptive
Dormand--Prince~4(5) scheme \citep{dormandprince:1980} from a static
observer tetrad, and disk emission follows a Novikov--Thorne thin disk
\citep{novikovthorne:1973,pagethorne:1974}, with images written to FITS.
The JP metric \citep{johannsen:2011} enters as a drop-in \texttt{Metric};
the integrator, camera, and renderer never branch on metric type. The
implementation passes machine-precision validation gates: metric--inverse
identity and Christoffel symmetry to $10^{-12}$; flat- and horizon-limit
recovery; closed circular photon orbits; null-condition drift bounds; and
convergence of the Schwarzschild shadow to $3\sqrt{3}\,M$ and of the Kerr
landmarks. We sweep $\epsilon_3\in[-2.5,+0.13]$ (16 values) at $a=0.9$,
inclinations $i\in\{17^\circ,60^\circ\}$, $512^2$ pixels. A ring extractor
bins each image into $180$ azimuthal sectors and records the
radially-outermost intensity peak per sector (integer-pixel radius),
reducing the image to the circularity $\delta_r/\langle r\rangle$ (the
ratio of the per-sector radius RMS to its mean). Code and data are
reproducible at repository tag \texttt{phase-3-complete}; the sweep was
generated at commit \texttt{4f711c0}.

\section{Results}
\emph{Photon-orbit bifurcation.} At $a=0.9$ the prograde equatorial
photon orbit ceases to exist above $\epsilon_3^{\rm crit}=0.1212$: beyond
this cusp the prograde shadow edge is set by the event horizon rather
than by a photon sphere. \emph{ISCO disappearance.} The prograde
innermost stable circular orbit disappears for $\epsilon_3\in(0.13,0.20)$;
there the Novikov--Thorne inner edge is undefined and the corresponding
sweep frames are skipped. Both are sharp, qualitative changes distinct
from the smooth-deformation regime and testable by future higher-precision
observations. \emph{Circularity trend.} Figure~\ref{fig:asym} shows
$\delta_r/\langle r\rangle$ versus $\epsilon_3$. On the monotone
sub-interval $\epsilon_3\in[-2.5,-0.2]$ at $i=17^\circ$ the circularity
rises smoothly from $0.121$ to $0.198$, plateauing near $0.195$ for
$|\epsilon_3|\lesssim0.1$. Figure~\ref{fig:gallery} shows the $i=17^\circ$
disk images rounding as $\epsilon_3$ increases. \emph{No EHT bound.} The
EHT \S7.4 circularity band ($0.055\pm0.005$) lies entirely below the
measured $\delta_r/\langle r\rangle$ ($\geq0.12$ throughout the sweep);
inverting the band against the monotone curve yields no two-sided
$\epsilon_3$ bound---the constraint runs off the low-$\epsilon_3$ edge.

\section{Discussion}
The $\sim\!0.12$ floor is not physical asymmetry but the noise floor of
integer-pixel peak finding: snapping each per-sector radius to the
nearest pixel ($\sim\!0.06\,M$ at $512^2$) caps the recoverable
circularity well above the EHT precision. Furthermore
$\delta_r/\langle r\rangle$ is an all-azimuthal-mode \emph{radius} RMS,
whereas the EHT \S7.4 measure is a \emph{diameter} spread in which a pure
$m=1$ asymmetry cancels; since the former bounds the latter from above,
our consistency test is conservative. Reaching the EHT's $\sim\!5.5\%$
circularity demands sub-pixel (e.g.\ parabolic) peak localization and a
diameter-matched estimator---the natural next step before $\epsilon_3$
can be constrained. The photon-orbit and ISCO predictions are insensitive
to extraction precision and stand independent of this limitation.

\begin{figure}[t]
\centering
\includegraphics[width=0.62\linewidth]{figures/asymmetry_vs_epsilon3.pdf}
\caption{Ring circularity $\delta_r/\langle r\rangle$ versus the JP
deviation $\epsilon_3$ at $a=0.9$ for $i=17^\circ$ and $60^\circ$. The
grey band is the EHT 2019 \S7.4 circularity ($0.055\pm0.005$); it lies
below both curves, so no two-sided $\epsilon_3$ bound follows. The dashed
line marks the photon-orbit cusp $\epsilon_3^{\rm crit}\approx0.12$ and
the shaded vertical band the ISCO-disappearance interval $(0.13,0.20)$.
\label{fig:asym}}
\end{figure}

\begin{figure}[t]
\centering
\includegraphics[width=\linewidth]{figures/image_gallery.pdf}
\caption{Novikov--Thorne disk images at $i=17^\circ$, $a=0.9$, for
$\epsilon_3\in\{-2.5,-1.0,0.0,+0.13\}$ on a shared logarithmic intensity
scale. The shadow rounds and contracts as $\epsilon_3$ increases toward
the photon-orbit cusp. \label{fig:gallery}}
\end{figure}

\begin{acknowledgments}
\software{nom-tam-fits, JUnit~5, NumPy \citep{numpy},
Matplotlib \citep{matplotlib}, Astropy \citep{astropy}, pandas}.
Computations ran on a MacBook~Pro~M3 (16~GB). This Note uses public EHT
data products.
\end{acknowledgments}

\bibliography{refs}
\bibliographystyle{aasjournal}

\end{document}
